% Paste-ready English introduction; decoupling-first revision on 2026-09-24.
% Flow: multi-sensor detection -> spatial/adaptive fusion -> why decouple
% -> object guidance for detection -> ObjDec -> evidence -> three contributions.
% Requires natbib (normally supplied by the paper template).
% Bibliography: references.bib
% This is a section fragment, not a standalone document.

\section{Introduction}
\label{sec:introduction}

Accurate 3D object detection is essential for autonomous driving.
Cameras, LiDAR, and 4D radar observe a scene through different sensing mechanisms, providing visual appearance, geometric structure, and radar-return information.
Combining these observations can compensate for limitations of individual sensors and provide a richer basis for detecting objects.
Realizing this potential requires integrating features with different sensing characteristics and spatial sampling patterns into a joint representation that supports detection, motivating continued research on multi-sensor fusion~\citep{liu2023bevfusion,park2025mome,yue2025roburcdet}.

Existing approaches have advanced multi-sensor detection through spatial interaction and adaptive fusion.
Establishing correspondence in a common spatial domain allows features from different sensors to exchange information about the same regions.
Attention and gating mechanisms further adjust their contributions, while foreground modeling and reliability estimation help account for differences in input quality.
Together, these developments provide an effective foundation for combining sensor observations under varying scene and sensing conditions~\citep{chae2024threedlrf,huang2025l4dr,chae2025dlrfusion,paek2025asf,park2026raf}.

Nevertheless, spatial correspondence and adaptive weighting do not by themselves distinguish the information that can be shared across sensors from the cues specific to each modality.
Even within aligned local regions, these two types of information remain mixed in the input features, although they play different roles in fusion: shared information can provide cross-modal support, while modality-specific cues can preserve complementary observations.
Learning consistency alone may suppress useful differences, whereas emphasizing differences alone may overlook evidence shared across sensors.
This motivates approaching fusion through representation decoupling: explicitly learning shared and modality-specific representations and using both to construct the fused detection features.
Such a formulation builds on advances in learning and retaining common and distinct information for multimodal prediction~\citep{liang2023factorcl,liu2025featurecausality,qian2026decalign,wang2026spfd}.

For dense 3D detection, this decoupling also needs a clear connection to detection targets.
Shared responses may originate from either objects or background, and modality-specific differences may reflect useful observations or noise.
Separating the two types of information therefore does not by itself ensure their relevance to detection.
We address this issue by grounding representation constraints in object regions and using object-related information to guide the contribution of the learned features to fusion.
Object guidance thus gives the learning and use of decoupled representations a detection-specific focus.

We introduce \textbf{ObjDec}, a multi-sensor 3D detection architecture centered on shared and modality-specific representation decoupling.
Within spatially corresponding bird's-eye-view (BEV) patches, separate mappings construct the two representations from each sensor token, and both participate in building the fused detection features.
Object-region supervision encourages cross-modal agreement in shared representations, preserves modality-specific differences, and limits similarity between the two branches.
The learned representations further determine feature updates, modality contributions, and object context, allowing their information to shape fusion inputs and cross-modal interaction.
Foreground gating and context modulation implement this connection between decoupled representations and fusion computation.
The representation mappings operate at every patch; object regions specify where the representation constraints are applied during training.
Ground-truth boxes provide training supervision only, while inference predicts all fusion-control signals from input features without prior detections or weather labels.

We evaluate ObjDec primarily on K-Radar.
At a confidence-filtering threshold of $0.3$ on v1, ObjDec achieves $88.36\%$ $\mathrm{AP}_{\mathrm{3D}}$ at IoU $0.3$ and $88.10\%$ $\mathrm{AP}_{\mathrm{BEV}}$ at IoU $0.5$.
Compared with the archived results associated with the released fusion baseline under the same filtering setting, these scores improve by $8.04$ and $7.77$ percentage points, respectively, with approximately $1.7\%$ more parameters.
Component ablations assess the contribution of each design, while representation visualizations and spatial gate analyses examine the learned feature organization and fusion behavior.
We also evaluate a decoupled-control variant on K-Radar v2 and train and test the architecture on V2X-Radar-V, examining its applicability beyond the primary category and sensor setting under the specified dataset protocols~\citep{paek2022kradar,paek2025asf,yang2025v2xradar}.

Our main contributions are as follows:
\begin{itemize}
    \item \textbf{A fusion architecture based on representation decoupling.}
    We develop a multi-sensor 3D detection architecture that explicitly learns shared and modality-specific representations in corresponding local regions and uses both to construct fused detection features.
    \item \textbf{Object-guided learning and use of decoupled representations.}
    We supervise representation decoupling within object regions and connect the learned representations to foreground gating, modality contribution control, and object context, guiding their use in detection-oriented fusion.
    \item \textbf{Evaluation of detection performance and model behavior.}
    We evaluate detection performance and component contributions on K-Radar, analyze representations and spatial gating, and examine the architecture under extended categories and target-dataset training on V2X-Radar-V.
\end{itemize}

\begin{figure}[p]
\centering
\fbox{\begin{minipage}[c][3cm][c]{0.9\linewidth}\centering\drafttodo{Insert the final author-approved motivation figure. This box only reserves its position.}\end{minipage}}
\caption{Motivation for object-guided shared and modality-specific representation learning and fusion.}
\label{fig:1}
\end{figure}

